%%  **********************************************
%
%   The lines beginning with "%" are comments. 
%   Theses comments explain how to use this template for
%   preparing the MCA Fifth Semester Mini-Project Report. 
%
%   The first non-comment line below contains 
%   the matter that go into the preamble of 
%   the input file. These are collected in a 
%   file named "preamble.tex". 
%
%   Do not make any changes in the first 
%   non-comment line below.
%
%%
\documentclass[11pt, oneside, a4paper]{book}
\usepackage{graphicx, fancyhdr, amsmath, times, enumerate, amssymb, calc, multirow, setspace, array}
\usepackage[a4paper,width=150mm,top=35mm,bottom=35mm,bindingoffset=6mm]{geometry}
%% Reference: https://www.overleaf.com/learn/latex/
%% How_to_Write_a_Thesis_in_LaTeX_(Part_2):_Page_Layout
%%
\newcommand{\VAtitle}[1]%
{\def\vtitle{#1}}%
\newcommand{\VAauthora}[1]%
{\def\vauthora{#1}}%
\newcommand{\VAauthorb}[1]%
{\def\vauthorb{#1}}%
\newcommand{\VAauthorc}[1]%
{\def\vauthorc{#1}}%
\newcommand{\VAadmissionyear}[1]%
{\def\vadmissionyear{#1}}%
\newcommand{\VAregisternumbera}[1]%
{\def\vregisternumbera{#1}}%
\newcommand{\VAregisternumberb}[1]%
{\def\vregisternumberb{#1}}%
\newcommand{\VAregisternumberc}[1]%
{\def\vregisternumberc{#1}}%
\newcommand{\VAguide}[1]%
{\def\vguide{#1}}%
\newcommand{\VAhod}[1]%
{\def\vhod{#1}}%
\newcommand{\VAdept}[1]%
{\def\vdept{#1}}%
\newcommand{\VAclass}[1]%
{\def\vclass{#1}}%
\newcommand{\VApaper}[1]%
{\def\vpaper{#1}}%
%%
\linespread{1.3}
\setcounter{secnumdepth}{4}
%%
%
%   No changes in the next non-comment line.
%
\begin{document}
%
%***************************************************
%
%   No changes in the next comment line.
%
\VApaper{20MCA245 Mini Project}%
%
%   In the next line, replace "Report Title" by 
%   the title of your project/thesis.
%
\VAtitle{Report Title}%
%
%   In the next line, replace "Student1" 
%   by name of first student. 
%   Write full name without, repeat without, Mr or Ms
%
\VAauthora{Student Name}%
%
s
%

%   In the next line, replace "RegNo1" by the
%   University Examination Register Number of first student.
%
\VAregisternumbera{VAS24MCA-20XX}% 
%
%   In the next line, replace "RegNo2" by the 
%   University Examination Register Number of second student.

%
%   In the next line, replace "2024" by the year of  
%   your admission to the College.
%
\VAadmissionyear{2024}
%
%   In the next line, replace "Guide" by 
%   the full name of your guide or supervisor 
%   with Mr. or Ms. or Dr. or Prof.
%
\VAguide{Ms. Siji K B}% 
%
%   No changes in the next three non-comment lines.
%
\VAhod{Dr. Reji C Joy}
%
\VAdept{Computer Applications}%
%
\VAclass{S3 MCA}%
%
%   The next line is code to print the first few pages: 
%       1. Title page
%       2. Bona Fide Certificate from MCA Dept
%       3. Declaration
%       4. Acknowledgment
%   No changes in the next non-comment line.
%
%
%   Redefining plain page style
%  
\fancypagestyle{plain}{%
\fancyhf{} % clear all header and footer fields
\fancyhead[L]{{\small \vtitle}}
\fancyhead[R]{\bf \thepage}
\fancyfoot[C]{\bfseries \thepage} 
\fancyfoot[L]%
{{\scriptsize Department of \vdept, 
Vidya Academy of Science \& Technology, Thrissur}}
\fancyfoot[R]%
{\includegraphics[width=0.75cm]{VidyaLogo}}%
\fancyfoot[C]{ }%
\renewcommand{\headrulewidth}{1pt}%
\renewcommand{\footrulewidth}{0.5pt}%
}%
%
\pagestyle{empty}
%
%%
\pagenumbering{roman}
\thispagestyle{empty}
%%
\quad
\begin{center}

\sffamily\small
%
\begin{spacing}{2}
{\fontsize{15.5}{20}\selectfont \bfseries \MakeUppercase{\vtitle}}\\[5cm]
\end{spacing}
%

\begin{spacing}{1.25}
A Mini Project Report submitted by\\[0.25cm]
{\bfseries \vauthora}

{(\vregisternumbera)}\\[0.1mm]


\quad\\[0.25cm]
to APJ Abdul Kalam Technological University\\
in partial fulfillment of the requirements for the award of the degree of\\
Master of Computer Applications
\end{spacing}
%
\quad\\[5cm]
%   Logo
%
\includegraphics[width=0.12\textwidth]%
{VidyaLogo}\\[0.3cm]
%
%   Department particulars
%
\begin{spacing}{1.25}
{\normalsize \sffamily \bfseries Department of \vdept}\\
{\small \sffamily Vidya Academy of Science \& Technology\\ 
\small Thalakkottukara, Thrissur - 680 501}\\
{\sffamily October 2025}
\end{spacing}
%%
\end{center}
%%
%\newpage
%%

%
\clearpage
%
\addcontentsline{toc}{chapter}{Certificate from College}
%
%\vspace*{\fill}
\quad\\[0.25cm]
\pagestyle{plain}
\begin{center}
%
%
%   Department particulars
%
{\Large \bf Department of Computer Applications \par  }%\\[0.2cm]
{\large \bf Vidya Academy of Science \& Technology\par}%\\
{\normalsize \bf Thalakkottukara, Thrissur - 680 501\\
({\tt http://www.vidyaacademy.ac.in})\par}
\qquad\\[0.5cm]
%
%   Logo
%
\includegraphics{VidyaLogo}\\[0.75cm]
%
{\Large \bf CERTIFICATE}\\[0.75cm]
%
\end{center}
%
%   Certificate statement
%
This is to certify that the report titled 
{\bf \vtitle} is a bona-fide record of the 
work related to the paper \vpaper\  done  by 

\begin{center}
{\bf \vauthora\  
(Reg. No. \vregisternumbera) } \\ 

\end{center}
   of  \vclass\  
(\vadmissionyear\ admissions) class 
of Vidya Academy of Science \& Technology, 
Thrissur - 680501 in partial fulfillment of the 
requirement for the award of the Degree of Master of \mbox{Computer} Applications  of APJ Abdul Kalam Technological University. \\
%
%   Spaces for signatures 
%
\begin{center}
\begin{tabular}{p{0.10\textwidth}p{0.35\textwidth}p{0.10\textwidth}p{0.35\textwidth}}
%
\multicolumn{2}{l}{\bf Guide/Supervisor}    
&\multicolumn{2}{l}{\bf Head of Department}   \\
%
Name & : \vguide    & Name & : \vhod \\
%
Signature&: ...........................................
\qquad\quad   & 
Signature&: ...........................................\\
%
Date &: \today    & Date &: \today \\
%
\end{tabular}
\\[0.75cm]
%
%   Space for seal of Department
%
{\small (Seal of Department of \vdept)}
\end{center} 
%
\vspace*{\fill}
%
%   End of Certificate
%  
\clearpage
%
\addcontentsline{toc}{chapter}{Declaration by Student}
\thispagestyle{empty}
\chapter*{Declaration}
%
I, {\bf \vauthora} , studying in Third Semester MCA (2024 admissions) class of Vidya Academy of Science \&\  Technology, Thrissur -- 680501,    hereby declare that the project report 
(``\vtitle'') submitted by me  is a bona fide work done by me
under supervision of \vguide. This submission represents my ideas in
my own words and where ideas or words of others have been included, I have adequately
and accurately cited and referenced the original sources. I also declare that I have
adhered to ethics of academic honesty and integrity and have not misrepresented or
fabricated any data or idea or fact or source in my submission. I understand that any
violation of the above will be a cause for disciplinary action by the institute and/or the
University and can also evoke penal action from the sources which have thus not been
properly cited or from whom proper permission has not been obtained. This report has
not been previously formed the basis for the award of any degree, diploma or similar title
of any other University. 

\qquad\\[1cm]
\begin{tabular}{p{0.05\textwidth}p{0.27\textwidth}p{0.01\textwidth}p{0.21\textwidth}p{0.35\textwidth}}
%
Place &: Thrissur -- 680501\qquad      &  &Signature of student  &: ......................................\\
Date  &: \today                        &   &Name of student       &: \vauthora \\
     
\end{tabular}
\clearpage
%
\addcontentsline{toc}{chapter}{Acknowledgment}
%
\chapter*{Acknowledgment}
%
\index{acknowledgment}
I wish to record my indebtedness and thankfulness 
to all who helped me prepare this Report\index{seminar report} titled 
`` \vtitle\ ''  and present it in a satisfactory way. This Report is part of my work related to the paper \vpaper.

I am especially thankful to my
guide and supervisor \vguide\  in the Department of \vdept\  
for giving me valuable suggestions and 
critical inputs in the preparation of this report. 
I am also thankful to \vhod, 
the Head of \vdept\  Department
for encouragement. 

My friends in my class have always been 
helpful and I am grateful to them for 
patiently listening to my presentations on my work related to the Project. 

\begin{flushright}
\vauthora\\ {\footnotesize (Reg. No. \vregisternumbera)}\\
\vclass\  (\vadmissionyear\  Admissions)\\
Vidya Academy of Science \& Technology\\
Thrissur -680 501.
\end{flushright}
%
\clearpage
%
%
%****************************************************
%
%   No changes in the next two non-comment lines. These
%   lines are intended to create the Synopsis page and 
%   to add the the item "Synopsis" to 
%   table of contents.
%
\chapter*{Synopsis}
\addcontentsline{toc}{chapter}{Synopsis}
%
%   Enter below a synopsis of the project report. 
%   The two lines given
%   are to be deleted before entering actual abstract.

This is only filler material. This is an outline, a template, 
for preparing mini project reports. The instructions given 
herein are to be followed verbatim.
%
%***************************************************
%
%   The next three lines creates a Table of Contents 
%   in a new page. 
%   Do not delete this. There must be a Table of Contents.
%
\clearpage
\addcontentsline{toc}{chapter}{Table of Contents}
\tableofcontents
%
%   The next three lines create a page containing a 
%   List of Tables. Delete them if there are no 
%   tables in your project report. 
%
\clearpage
\addcontentsline{toc}{chapter}{List of Tables}
\listoftables
%
%   The next three line create a page containing a 
%   List of Figures. Delete them if there are no 
%   figures in your project report.
%
\clearpage
\addcontentsline{toc}{chapter}{List of Figures}
\listoffigures
%
%***************************************************
%
\mainmatter
%
\clearpage
\quad\vfill
\begin{center}
{\Huge \bf PROJECT REPORT}
\end{center}
\vfill
\clearpage
%
%
%   The main contents of the paper begin here. 
% 
%
%   *********************  CHAPTER ***********************
%
%  
\chapter{INTRODUCTION}

%
\section{Problem Definition}
%
\section{Objectives}
%
\section{Purpose of the Project}

%
%   *********************  CHAPTER ***********************
%
%
\chapter{SYSTEM ANALYSIS}
%
\section{The Existing System}	

%< Describe your various components here >
\section{Proposed System}
%
%
%   *********************  CHAPTER ***********************
%
%
\chapter{FEASIBILITY STUDY}
%
\section{Technical Feasibility}
%
\section{Economical Feasibility}
%
\section{Operational Feasibility}
%
\chapter{SYSTEM DESIGN}
%
\section{Input Design}
%
\section{Output Design}
%
\section{Database Design}
%
%
%   *********************  CHAPTER ***********************
%
%
\chapter{IMPLEMENTATION, TESTING AND VALIDATION}
%
\section{Implementation}
%
\section{Testing Methods}
% < Describe the various Testing methods used here >
\section{Validation}
%	< Validation done in the project >
%
%   *********************  CHAPTER ***********************
%
\chapter{CONCLUSION}
%
%   *********************  CHAPTER ***********************
%
\chapter{SCOPE FOR FUTURE ENHANCEMENT}
%
%
%   *********************  CHAPTER ***********************
%
%%%%%%%%%%%%%%%%%%%%%%%%%  APPENDICES %%%%%%%%%%%%%%%%%%%%%
%
\clearpage
\addcontentsline{toc}{chapter}{APPENDICES}
\appendix
\quad\vfill
\begin{center}
{\Huge \bf APPENDICES}
\end{center}
\vfill
\clearpage
%
%
%   *********************  CHAPTER ***********************
%
%
\chapter{SCRUM PROCESS ARTIFACTS}
%
\section{Product Backlog}
%
% The following table is only a sample. 
%
\renewcommand{\arraystretch}{1.25}
\begin{center}
\begin{tabular}{|m{0.08\textwidth}|m{0.15\textwidth}|m{0.5\textwidth}|m{0.1\textwidth}|}
\hline
{\bf Priority} & {\bf Product backlog items} & {\bf User story} & {\bf Estimate (Hours)}\\
\hline
1 & Database creation & As an operations engineer, I want to be able to store all customer information.  & 240 \\
\hline
2 & Login page & As a site member I want to login to the site. & 160 \\
\hline
3 & Category page & As a site member, I want to be able to look for different categories of brands. & 400\\
\hline
4 & Payment process & As a site member, I want to be able to make payments. & 240\\
\hline
5 & Contact page & As a site member, I want to be able to find contact information of the site. & 80 \\
\hline
6 & Banner area & As a marketing personnel, I want to be able to make advertisement. & 40\\
\hline
\end{tabular}
\end{center}
%
\section{Scrum Meeting Details}
%
\begin{center}
\begin{tabular}{|m{0.03\textwidth}|m{0.05\textwidth}|m{0.3\textwidth}|m{0.2\textwidth}|m{0.2\textwidth}|}
\hline
Sl. No. & Date & Sprint backlog & Product shippable Increment & Important Decisions \\
\hline
& & & &  \\
\hline
\end{tabular}
\end{center}

\renewcommand{\arraystretch}{1}
%
%   *********************  APPENDIX ***********************
%
\chapter{UML DIAGRAMS}
\section{Class Diagram}
%
\begin{figure}
   \centering
    \includegraphics[width=0.5\linewidth]{chatgptaopenai.jpg}
    \caption{Enter Caption}
    \label{fig:enter-label}
\end{figure}
\begin{figure}
   \centering
    \includegraphics[width=0.5\linewidth]{chatgptaopenai.jpg}
    \caption{Chatgpt}
    \label{fig:enter-label}
\end{figure}
\section{State Diagram}
%
%
%   *********************  APPENDIX ***********************
%
\chapter{TABLE STRUCTURE}
\begin{table}
 \centering
  \begin{tabular}{|l|l|l|l|l|}
    \hline 
   Sl. No.&  Field&  Datatype&  Constraint& Description \\
     \hline  
    1 & Stud\_No & Number &Primary Key  & Student Register Number\\ 
      \hline 
     2 & Stud\_Name  & Text & Not Null & Student Name\\
     \hline 
     3  & Phone\_Number & Numeric & Not Null & Phone number \\ 
				\hline
    \end{tabular}
    \caption{Student table}
\label{table:Student}
\end{table}
\begin{table}
 \centering
  \begin{tabular}{|l|l|l|l|l|}
    \hline 
   Sl. No.&  Field&  Datatype&  Constraint& Description \\
     \hline  
    1 & Stud\_No & Number &Primary Key  & Student Register Number\\ 
      \hline 
     2 & Stud\_Name  & Text & Not Null & Student Name\\
     \hline 
     3  & Phone\_Number & Numeric & Not Null & Phone number \\ 
				\hline
    \end{tabular}
    \caption{Master table}
\label{table:Student}
\end{table}

%   *********************  APPENDIX ***********************
%
\chapter{SAMPLE SCREENSHOTS}
%
%
% ***   Creating Bibliography   ***
%
% Do not delete the following two lines
%
\clearpage
\addcontentsline{toc}{chapter}{Bibliography}
%
%   The markups for creating the bibliography 
%   begins here. Do not change the two lines 
%   "\begin{thebibliography}{99}" and 
%   "\end{thebibliography}".
%   One sample bibliography item is included. 
%   Each new item is to be preceded by
%   "\bibitem{itemx}".
%   Add additional items if there are any.
%   Bibliography styles:
%   Titles of books   : Italics  {use {\em Title} )
%   URLs of websites  : Type-writer (use {\tt Title} )
%   Titles of papers  : In double quotes (use ``Title")
%
\begin{thebibliography}{99}
%
\bibitem{item1}
{\em First Lessons in \LaTeX}, by Dr. V N Krishnachandran, 
Vidya Academy of Science \& Technology, 
Thrissur - 680 501, 2011.
%
\end{thebibliography}
%
%
%   Printing the last page
%
%
\newpage
\thispagestyle{empty}
\vspace*{\fill}
\begin{flushright}
\includegraphics{VidyaLogo}\\[0.5cm]
{\Large \bf \sf  Department of \vdept\ }\\
{\sf Vidya Academy of Science \& Technology\\
Thalakkottukara, Thrissur - 680 501\\
({\tt http://www.vidyaacademy.ac.in})}
\end{flushright}
%
%
%
%   ***   The end   ***
%
%
\end{document}